% Template-wise test suite for the RVV 1.0 vector loads and stores on the Kites simulator.
% Self-contained: compile with  pdflatex vector_ldst_test_suite.tex
% Tests: examples/vector_ldst_suite/run_suite.py (programs in examples/vector_ldst_suite/programs/)
\documentclass[11pt,a4paper]{article}

\usepackage[margin=1in]{geometry}
\usepackage[T1]{fontenc}
\usepackage{lmodern}
\usepackage{amsmath}
\usepackage{array}
\usepackage{booktabs}
\usepackage{enumitem}
\usepackage{xcolor}
\usepackage{listings}
\usepackage[hidelinks]{hyperref}

\lstset{
    basicstyle=\ttfamily\small,
    columns=fullflexible,
    keepspaces=true,
    frame=single,
    breaklines=true,
    commentstyle=\color{gray},
    morecomment=[l]{\#}
}

\title{Test Suite for the Vector Load and Store Instructions}
\author{RISC-V Vector Extension (RVV 1.0) project}
\date{}

\begin{document}
\maketitle

\section*{How the tests work}

The suite contains 100 test cases, grouped by the instruction template they exercise. Each case is
a small assembly program, generated by \texttt{examples/vector\_ldst\_suite/run\_suite.py} and
saved in \texttt{examples/vector\_ldst\_suite/programs/<template>/}, so that it can also be opened
and stepped through in the GUI. Every program has the same four parts.

\begin{enumerate}
    \item Registers \texttt{v8}--\texttt{v31} are filled with a known byte pattern, \texttt{v0}
    with the mask, and \texttt{v16} onwards with the index values (indexed cases only).
    \item \texttt{vsetvli} sets the SEW, LMUL and AVL of the case with the tail-undisturbed and
    mask-undisturbed policies (\texttt{tu, mu}).
    \item The instruction under test is executed once.
    \item The new \texttt{vl} is read back with \texttt{vmset.m} and \texttt{vcpop.m}, so the check does not
    depend on CSR instructions. Then
    \texttt{v8}--\texttt{v31}, \texttt{vl} and the 256-byte store area are written to memory.
\end{enumerate}

The script runs every program with \texttt{kites-cli} on the single-cycle VM (\texttt{rvss}) and
on both pipelined VMs with hazard detection (\texttt{rv5-h-nf}, \texttt{rv5-h-f}). It compares the
dumped bytes with the result of a small reference model of the RVV~1.0 load and store rules,
written separately from the simulator's C++ code. The model covers the address of every element
for each template, register groups (EMUL), segment fields, masking, tail elements, the lengths of
whole-register and mask accesses, unsigned index offsets, 64-bit address wrap-around and
fault-only-first trimming. Every register byte, \texttt{vl} and every byte of the store area must
match.

For the error cases the run must stop with a specific message, and the suite also checks where
the error was detected. Most illegal forms are caught by the processor when the instruction
executes (``Run failed: Illegal instruction \ldots''), because they depend on \texttt{vtype}. Two
forms are already rejected by the assembler, because they are illegal for any \texttt{vtype}.
Access faults must name the faulting element.

\begin{table}[h]
\centering
\begin{tabular}{|l|c|c|c|}
    \hline
    \textbf{Template} & \textbf{Cases} & \textbf{Runs (3 VMs)} & \textbf{Passed} \\
    \hline
    Unit-stride                       & 24 & 72  & 72 \\
    Strided                           & 14 & 42  & 42 \\
    Indexed                           & 16 & 48  & 48 \\
    Segment                           & 15 & 45  & 45 \\
    Whole-register                    & 13 & 39  & 39 \\
    Mask                              & 9  & 27  & 27 \\
    Fault-only-first and access faults & 9 & 27  & 27 \\
    \hline
    \textbf{Total}                    & \textbf{100} & \textbf{300} & \textbf{300} \\
    \hline
\end{tabular}
\caption{Results of the load/store test suite.}
\label{tab:suite}
\end{table}

\paragraph{Negative controls.} To confirm that the checks can detect a wrong result, seven
mutated programs were checked against the expectations of the original case. All seven were
reported as failures: a masked load with its mask removed, a zero stride replaced by 4, a
three-field segment load replaced by a two-field one, \texttt{vl4re32.v} replaced by
\texttt{vl2re32.v}, \texttt{vlm.v} replaced by \texttt{vle8.v}, a fault-only-first load replaced by
an ordinary load, and an index EEW of 32 replaced by 8.

\paragraph{Unit tests.} The same behaviour is also covered at the C++ level by the 30 GoogleTest
tests in \texttt{test/vector\_load\_store\_test.cpp}. These tests also check the encodings
against the official RISC-V opcode tables, the assembler and disassembler round trip, execution
from a non-zero \texttt{vstart}, and undo and redo.

%==============================================================================================
\subsection*{Unit-Stride Loads and Stores (\texttt{vle<eew>.v}, \texttt{vse<eew>.v})}
\label{subsec:ldst-unit}

The following edge cases are covered for the unit-stride loads and stores, at element widths
8, 16, 32 and 64.

\begin{itemize}
    \item \textbf{Element widths and register groups}
    \begin{itemize}
        \item Every element width with SEW equal to EEW
        \item A full LMUL $= 8$ group with \texttt{vl} $=$ VLMAX $= 32$ (\texttt{v8}--\texttt{v15})
        \item EEW smaller than SEW (\texttt{vle8.v} at SEW $= 32$, EMUL $= 1/4$)
        \item EEW larger than SEW (\texttt{vle64.v} at SEW $= 16$, EMUL $= 4$)
        \item A fractional LMUL (\texttt{vle16.v} at SEW $= 32$, LMUL $= 1/2$)
    \end{itemize}

    \item \textbf{Vector length and tail elements}
    \begin{itemize}
        \item \texttt{vl} $= 0$ for a load and for a store, which change nothing
        \item \texttt{vl} $= 1$
        \item A partial LMUL $= 4$ group (\texttt{vl} $= 13$), with the tail running across
              register boundaries
        \item A store with \texttt{vl} $= 7$, where the memory after the last element is untouched
    \end{itemize}

    \item \textbf{Masking}
    \begin{itemize}
        \item A masked load and a masked store with the mask \texttt{0101\ldots}
        \item A masked load with an all-zero mask, which changes nothing
    \end{itemize}

    \item \textbf{Alignment}
    \begin{itemize}
        \item A 32-bit load from an address that is not a multiple of 4
        \item A 64-bit store to an address that is not a multiple of 8
    \end{itemize}

    \item \textbf{Illegal instruction conditions}
    \begin{itemize}
        \item EMUL larger than 8 (\texttt{vle64.v} at SEW $= 8$, LMUL $= 2$)
        \item A misaligned destination group (\texttt{v9} at LMUL $= 2$)
        \item A masked load whose destination is \texttt{v0} (rejected by the assembler)
        \item Any load after a \texttt{vsetvli} that set \texttt{vill}
    \end{itemize}
\end{itemize}

%==============================================================================================
\subsection*{Strided Loads and Stores (\texttt{vlse<eew>.v}, \texttt{vsse<eew>.v})}
\label{subsec:ldst-strided}

The following edge cases are covered for the strided loads and stores, where the byte stride
is taken from \texttt{rs2}.

\begin{itemize}
    \item \textbf{Stride values}
    \begin{itemize}
        \item A stride equal to the element size, which behaves like a unit-stride access
        \item A stride larger than the element size (every third word, and a stride of 128)
        \item A negative stride, which reads an array in reverse order
        \item A zero stride on a load, which copies one element into every lane
        \item A zero stride on a store, where every element goes to the same address and the
              last element wins
        \item A stride that is not a multiple of the element size (6 bytes for 32-bit
              elements), so the elements are misaligned
    \end{itemize}

    \item \textbf{Element widths}
    \begin{itemize}
        \item SEW $= 64$ with LMUL $= 2$
        \item EEW smaller than SEW (\texttt{vlse8.v} at SEW $= 32$, stride 3)
    \end{itemize}

    \item \textbf{Stores}
    \begin{itemize}
        \item A store with a positive stride, where the gaps between the elements are untouched
        \item A store with a negative stride
    \end{itemize}

    \item \textbf{Masking and vector length}
    \begin{itemize}
        \item A masked load and a masked store with the mask \texttt{0101\ldots}
        \item \texttt{vl} $= 0$
    \end{itemize}
\end{itemize}

%==============================================================================================
\subsection*{Indexed Loads and Stores (\texttt{vluxei}, \texttt{vloxei}, \texttt{vsuxei},
\texttt{vsoxei})}
\label{subsec:ldst-indexed}

The following edge cases are covered for the unordered and ordered indexed loads (gather) and
stores (scatter). The index register holds byte offsets with the instruction's index EEW, and
the data elements use SEW.

\begin{itemize}
    \item \textbf{Index widths}
    \begin{itemize}
        \item Index EEW 8, 16, 32 and 64 with 32-bit data, so that the index EMUL is $1/4$, $1/2$,
              $1$ and $2$
        \item 8-bit indices with 64-bit data at LMUL $= 2$
    \end{itemize}

    \item \textbf{Offset values}
    \begin{itemize}
        \item A permutation of the offsets, for the unordered and the ordered load
        \item An 8-bit offset of 252, which must be read as unsigned (not as $-4$)
        \item The same offset repeated in several lanes
        \item A 64-bit offset of $2^{64} - 4$, which wraps around to the base address minus 4
    \end{itemize}

    \item \textbf{Stores}
    \begin{itemize}
        \item An unordered scatter
        \item An ordered scatter with duplicate offsets, where the last element wins
        \item An ordered scatter with 8-bit indices at SEW $= 16$
    \end{itemize}

    \item \textbf{Masking and vector length}
    \begin{itemize}
        \item A masked gather and a masked scatter with the mask \texttt{0101\ldots}
        \item \texttt{vl} $= 0$
    \end{itemize}

    \item \textbf{Register overlap}
    \begin{itemize}
        \item The destination equal to the index register with the same EEW, which is legal
        \item The destination equal to the index register with a different EEW, which is illegal
    \end{itemize}
\end{itemize}

%==============================================================================================
\subsection*{Segment Loads and Stores (\texttt{vlseg}, \texttt{vlsseg}, \texttt{vl[u/o]xseg} and
the stores)}
\label{subsec:ldst-segment}

The following edge cases are covered for the segment loads and stores, which move records of
NFIELDS consecutive fields. Field $f$ of every element goes to register group
\texttt{vd} $+ f \cdot$ EMUL.

\begin{itemize}
    \item \textbf{Number of fields}
    \begin{itemize}
        \item Two fields: an array of $(x, y)$ pairs split into \texttt{v8} and \texttt{v9}
        \item Three fields: RGB bytes split into three registers
        \item Eight fields, the maximum
        \item Four fields at LMUL $= 2$, where NFIELDS $\times$ EMUL $= 8$ (the limit)
        \item EEW smaller than SEW (\texttt{vlseg2e8.v} at SEW $= 32$)
    \end{itemize}

    \item \textbf{Addressing modes}
    \begin{itemize}
        \item Unit-stride segments (\texttt{vlseg}, \texttt{vsseg})
        \item Strided segments, reading the first two fields of 12-byte records
              (\texttt{vlsseg2e16.v}) and writing with a 16-byte stride (\texttt{vssseg2e32.v})
        \item Indexed segments (\texttt{vluxseg2ei8.v}, \texttt{vsoxseg3ei16.v})
    \end{itemize}

    \item \textbf{Stores}
    \begin{itemize}
        \item Two registers interleaved back into pairs
        \item A masked three-field store
    \end{itemize}

    \item \textbf{Masking, tail and vector length}
    \begin{itemize}
        \item \texttt{vl} $= 3$ with the mask \texttt{0101\ldots}, so whole segments are skipped
              and the tail of every field register is kept
        \item \texttt{vl} $= 0$
    \end{itemize}

    \item \textbf{Illegal instruction conditions}
    \begin{itemize}
        \item NFIELDS $\times$ EMUL larger than 8 (three fields at LMUL $= 4$)
        \item Fields running past \texttt{v31} (four fields starting at \texttt{v30})
    \end{itemize}
\end{itemize}

%==============================================================================================
\subsection*{Whole-Register Loads and Stores (\texttt{vl<nf>re<eew>.v}, \texttt{vs<nf>r.v})}
\label{subsec:ldst-whole}

The following edge cases are covered for the whole-register loads and stores, which always move
NFIELDS complete registers ($\mathrm{NFIELDS} \cdot \mathrm{VLEN}/8$ bytes).

\begin{itemize}
    \item \textbf{Register counts}
    \begin{itemize}
        \item Loads of 1, 2, 4 and 8 registers (\texttt{vl1re8.v}, \texttt{vl2re16.v},
              \texttt{vl4re32.v}, \texttt{vl8re64.v})
        \item Stores of 1, 2, 4 and 8 registers (\texttt{vs1r.v} to \texttt{vs8r.v})
        \item The alias \texttt{vl1r.v} of \texttt{vl1re8.v}
    \end{itemize}

    \item \textbf{Independence from \texttt{vl} and \texttt{vtype}}
    \begin{itemize}
        \item \texttt{vl} $= 0$ for a load and for a store, which still move whole registers
        \item SEW $= 64$, LMUL $= 8$, \texttt{vl} $= 1$ and an all-zero mask, which do not
              change the result
        \item A load after a \texttt{vsetvli} that set \texttt{vill}, which still executes
    \end{itemize}

    \item \textbf{Alignment}
    \begin{itemize}
        \item A load from an address that is not a multiple of 4
        \item A two-register load into the odd register \texttt{v9} (rejected by the assembler)
    \end{itemize}
\end{itemize}

%==============================================================================================
\subsection*{Mask Loads and Stores (\texttt{vlm.v}, \texttt{vsm.v})}
\label{subsec:ldst-mask}

The following edge cases are covered for the mask loads and stores, which move
$\lceil \texttt{vl}/8 \rceil$ bytes whatever the SEW and LMUL are.

\begin{itemize}
    \item \textbf{Byte count}
    \begin{itemize}
        \item \texttt{vl} $= 10$ and \texttt{vl} $= 16$, which both move 2 bytes
        \item \texttt{vl} $= 17$, which moves 3 bytes
        \item \texttt{vl} $= 1$ on a store, which moves 1 byte
        \item \texttt{vl} $= 0$ on a load and on a store, which move nothing
        \item SEW $= 32$, LMUL $= 8$, \texttt{vl} $= 32$, which moves 4 bytes
    \end{itemize}

    \item \textbf{Register contents}
    \begin{itemize}
        \item The bytes of the destination register after the last loaded byte are kept
        \item The alias \texttt{vle1.v} of \texttt{vlm.v}
    \end{itemize}
\end{itemize}

%==============================================================================================
\subsection*{Fault-Only-First Loads and Access Faults (\texttt{vle<eew>ff.v},
\texttt{vlseg<nf>e<eew>ff.v})}
\label{subsec:ldst-fof}

The following edge cases are covered for the fault-only-first loads and for access faults in
the other templates. The tests use addresses at the top of the 64-bit address space, where the
last valid byte is $2^{64} - 2$.

\begin{itemize}
    \item \textbf{Fault-only-first loads}
    \begin{itemize}
        \item No fault, where the load behaves like \texttt{vle32.v} and \texttt{vl} is unchanged
        \item A fault at element 3 of a byte load, which sets \texttt{vl} $= 3$ and keeps the
              elements from 3 onwards unchanged
        \item A fault at element 2 of a 32-bit load, which sets \texttt{vl} $= 2$
        \item A fault in the second field of element 2 of a segment load, which sets
              \texttt{vl} $= 2$ and does not write any field of element 2
        \item A fault at element 0, which traps like an ordinary load
    \end{itemize}

    \item \textbf{Access faults in the other templates}
    \begin{itemize}
        \item A unit-stride load and a unit-stride store whose second element is out of range
        \item A strided load whose second element is out of range
        \item An indexed load whose second offset reaches address $2^{64} - 1$
        \item In every case the message names the instruction, the faulting element and its
              address
    \end{itemize}
\end{itemize}

\end{document}
